\documentclass[a4paper,12pt]{article}
\usepackage[utf8]{inputenc}
\usepackage[T1]{fontenc}
\usepackage[ngerman]{babel}
\usepackage{amsmath, amssymb}
\usepackage{tikz}
\usetikzlibrary{shapes.geometric, arrows.meta, positioning, fit, backgrounds}
\usepackage{geometry}
\geometry{left=3cm, right=3cm, top=2.5cm, bottom=2.5cm}
\usepackage{enumitem}
\usepackage{parskip}

\tikzset{
  box/.style={rectangle, draw, rounded corners=3pt, text width=4.5cm,
              align=center, minimum height=1.0cm, font=\small, inner sep=4pt},
  arrow/.style={-{Latex[length=2mm]}, thick},
}

\begin{document}

\begin{center}
  {\LARGE\bfseries Studierhinweise zu Lektion N}\\[0.5em]
  {\large MODULNAME (Modul NUMMER)}
\end{center}

\vspace{1em}

% Einleitungstext: 1--3 Absätze über Inhalt und Aufbau der Lektion.

% -----------------------------------------------------------------------
\section*{Struktur der Lektion N}

\begin{center}
\begin{tikzpicture}[node distance=0.6cm and 0.8cm]

  % Spaltenköpfe
  \node[font=\bfseries\small] (h1) at (0,0)   {Kapitel 1: TITEL};
  \node[font=\bfseries\small] (h2) at (6.5,0) {Kapitel 2: TITEL};

  % Boxen Kapitel 1
  \node[box, below=0.8cm of h1] (a1) {1.1 ABSCHNITT\\Stichworte};
  \node[box, below=0.5cm of a1] (a2) {1.2 ABSCHNITT\\Stichworte};

  % Boxen Kapitel 2
  \node[box, below=0.8cm of h2] (b1) {2.1 ABSCHNITT\\Stichworte};
  \node[box, below=0.5cm of b1] (b2) {2.2 ABSCHNITT\\Stichworte};

  % Pfeile innerhalb Kapitel
  \draw[arrow] (a1) -- (a2);
  \draw[arrow] (b1) -- (b2);

  % Kapitelübergreifende Pfeile (nur wenn nötig)
  \draw[arrow] (a2.east) -- ++(0.3,0) |- (b1.west);

\end{tikzpicture}
\end{center}

\newpage

% -----------------------------------------------------------------------
\section*{Zielelement 1.1 -- TITEL DES ABSCHNITTS}

\subsection*{Lerninhalte}

\begin{center}
\begin{tikzpicture}[node distance=0.6cm and 1.0cm]
  \node[box] (k1) {Begriff/Konzept 1};
  \node[box, below=0.5cm of k1] (k2) {Begriff/Konzept 2};
  \node[box, below=0.5cm of k2] (k3) {Begriff/Konzept 3};
  \draw[arrow] (k1) -- (k2);
  \draw[arrow] (k2) -- (k3);
\end{tikzpicture}
\end{center}

\subsection*{Lernziele}

Nach Durcharbeiten dieses Abschnitts sollten Sie:
\begin{itemize}
  \item ...können,
  \item ...kennen,
  \item ...verstehen und anwenden können.
\end{itemize}

\subsection*{Selbstkontrollelement 1.1}

% Eine kurze, selbst beantwortbare Aufgabe zum Kernbegriff des Abschnitts.

\newpage

% -----------------------------------------------------------------------
% Weitere Zielelemente nach demselben Schema...

\end{document}
